% Lösungen Daten (zusammenbau v0.4, Heft)
\einheitenkopf{Daten}

\erg{1}{a) $26$ Kinder (alle) \quad b) Rot $= 11$, Blau $= 7$ \quad c) Anzahlen $= 1$, $2$, $3$, $4$ (für $1$, $2$, $3$, $5$)}
\erg{2}{a) $\frac{17}{26} \approx 0{,}654 \approx 65{,}4\,\%$ \quad b) $\frac{16}{28} = \frac{4}{7} \approx 0{,}571 \approx 57{,}1\,\%$}
\erg{3}{a) $20$ je Linie ($100 : 5$) \quad b) $40$ Besucher \quad c) $65$ Punkte ($60$ und eine Hilfslinie zu je $5$)}
\erg{4}{a) $270$ Tausend $= 270\,000$ Besucher \quad b) Dienstag, Freitag und Sonntag; der Samstag mit genau $220$ zählt nicht \quad c) $400$ Tausend hat genau das Jahr 2022, es zählt nicht; weniger hatten die Jahre 2019 und 2021 \quad d) Balken Mi bis $230$ (drei Hilfslinien über $200$) \quad e) Säule So bis $1{,}45$ (die Hilfslinie zwischen $1{,}4$ und $1{,}5$)}
\erg{5}{a) a) $38 : 3{,}8 = 10$ Mio. je Kästchen; b) $5{,}9 \cdot 10 = 59$, also Italien; c) Frankreich $= 68 : 10 = 6{,}8$ Kästchen hoch \quad b) ausgelassen \quad c) a) $84 : 7 = 12$ je Linie; b) A $= 72$ Kirsche, B und E $= 84$ Apfel und Birne, C $= 48$ Banane, D $= 36$ Traube \quad d) a) $180 : 6 = 30$ je Linie; b) A $= 90$ Fisch, B und D $= 60$ Hamster und Vogel, C $= 180$ Katze, E $= 150$ Hund}
\erg{6}{a) Winkel in Grad \quad b) $3$ cm ($30$ mm) \quad c) Abschnitte $= 4$ cm, $3$ cm, $3$ cm, beschriftet}
\erg{7}{a) Abschnitte $= 2$ cm, $4{,}5$ cm, $1$ cm, $2{,}5$ cm, beschriftet \quad b) $100 - 21 - 58 - 7 - 2 = 12\,\%$ \quad c) $9 : 76 \cdot 360^\circ = 42{,}6\ldots^\circ \approx 43^\circ$ \quad d) größter Sektor täglich, dann selten, dann mehrmals; Winkel „nie“ $= 0{,}08 \cdot 360^\circ = 28{,}8^\circ$}
\erg{8}{a) a) Kühlen ist der drittgrößte Sektor ($540$ kWh, $54^\circ$), zwischen Warmwasser und den beiden gleich großen Sektoren Kochen und Sonstiges; b) $180 : 3\,600 \cdot 360^\circ = 18^\circ$ \quad b) a) Wäsche ist der drittgrößte Sektor ($63$ Liter, $54^\circ$), zwischen Toilette und Küche; b) $14 : 420 \cdot 360^\circ = 12^\circ$ \quad c) a) zu Fuß, der kleinste Anteil; b) $0{,}124 \cdot 360^\circ = 44{,}64^\circ \approx 45^\circ$; c) Sektor $45^\circ$ am vorhandenen Radius angetragen \quad d) a) Biomasse, der kleinste Anteil; b) $0{,}183 \cdot 360^\circ = 65{,}88^\circ \approx 66^\circ$; c) Sektor $66^\circ$ am vorhandenen Radius angetragen \quad e) a) $5$ von $13$: $5 : 13 = 0{,}3846\ldots \approx 38{,}5\,\%$; b) Winkel $0{,}3846 \cdot 360^\circ \approx 138^\circ$, Sektor ab dem Radius \quad f) a) $5$ von $9$: $5 : 9 = 0{,}5555\ldots \approx 55{,}6\,\%$; b) Winkel $0{,}5556 \cdot 360^\circ = 200^\circ$, Sektor ab dem Radius}
\erg{9}{a) Spannweite \quad b) Minimum $= 9$, Maximum $= 22$ \quad c) Er steigt, weil der neue Wert über dem bisherigen Mittelwert liegt.}
\erg{10}{a) Minimum $= 162{,}4$ ct, Maximum $= 171{,}9$ ct \quad b) Minimum $= 95$, Maximum $= 240$, Spannweite $= 145$ Mio. \quad c) $4{,}62 - 4{,}18 = 0{,}44$ m \quad d) $158{,}6 - 149{,}8 = 8{,}8$ ct \quad e) $1{,}73 - 1{,}61 = 0{,}12$ € \quad f) sortiert: $2$, $4$, $6$, $7$, $9$, $9$, $11$ → Median $= 7$ °C}
\erg{11}{a) sortiert: $9$, $13$, $14$, $16$, $19$, $19$ → Median $= (14 + 16) : 2 = 15$ °C \quad b) $(6 + 30 + 45) : 3 = 81 : 3 = 27$ \quad c) $5{,}8 : 4 = 1{,}45$ m \quad d) $20{,}4 : 4 = 5{,}10$ m \quad e) Spannweite $= 8{,}3 - 6{,}2 = 2{,}1$ Mio.; Mittel $= 59{,}6 : 8 = 7{,}45$ Mio.}
\erg{12}{a) $(120 + 90 + 160 + 150 + 140 + 210 + 250) : 7 = 1\,120 : 7 = 160$ \quad b) Minimum $= 310$ (2023), Maximum $= 390$ (2021); Mittel $= (350 + 390 + 340 + 310 + 360) : 5 = 350$ \quad c) größter $= 92$ (Juli), kleinster $= 16$ (Februar); Spannweite $= 92 - 16 = 76$ mm \quad d) $412\,879 : 19 = 21\,730{,}47\ldots \approx 21\,730$ Zuschauer (oder etwa $21\,700$) \quad e) $634{,}8 : 12 = 52{,}9$ mm – die Behauptung stimmt \quad f) Summe $= 7 \cdot 18 = 126$; Samstag $= 126 - 106 = 20$ °C}
\erg{13}{a) Aussage 2 – der häufigste Wert ist $= 3$ (viermal), nicht $5$ (dreimal); Aussage 1 stimmt: $6 - 2 = 4$ \quad b) Aussage 1: $250 - 90 = 160$, wahr; Aussage 2: ein Viertel von $160$ ist $= 40$, $160 + 40 = 200$, nicht $250$ – falsch (es sind mehr als die Hälfte mehr) \quad c) $(1{,}62 + 1{,}64 + 1{,}61 + 1{,}66 + 1{,}71) : 5 = 8{,}24 : 5 = 1{,}648$ € – das ist unter $1{,}65$, die Behauptung stimmt (nur die fünf Werktage mitteln) \quad d) a) $440 : 11 = 40$ Jahre; b) $61 + 19 = 80 = 2 \cdot 40$: die beiden haben zusammen genau den Durchschnitt, also bleibt er $(440 - 80) : 9 = 40$ \quad e) a) $297 : 11 = 27$ Jahre; b) $35 + 19 = 54 = 2 \cdot 27$: die beiden haben zusammen genau den Durchschnitt, also bleibt er $(297 - 54) : 9 = 27$}
\erg{14}{a) bei null \quad b) Nein – im März waren es genau $= 60$, „mehr als“ heißt darüber. \quad c) ausgelassen}
\erg{15}{a) Teil 1 falsch: 2021 ist die Zahl gestiegen ($348$ auf $361$). Teil 2 falsch: von $348$ auf $296$ ist ein Rückgang um $= 52$, weit unter der Hälfte.}
\erg{16}{a) ausgelassen \quad b) ausgelassen \quad c) Die Achse beginnt bei $2{,}10$ statt bei $0$; die Säulen zeigen nur den Teil darüber: Do $0{,}04$, Fr $0{,}08$. Dieser Teil verdoppelt sich, der Preis steigt aber nur um $= 0{,}04$ €, also knapp $2\,\%$. \quad d) Die Achse beginnt bei $18$ statt bei $0$; sichtbar sind nur $1$ und $3$ Grad über dem Achsenanfang. Diese Teile verhalten sich wie $1 : 3$, die Temperatur stieg aber nur um $= 2$ Grad, gut $10\,\%$. \quad e) Aussage 1 falsch: $52\,\%$ sind mehr als die Hälfte. Aussage 2 richtig: jede 25. ist $\frac{1}{25} = 4\,\%$, wie im Diagramm; die Zahl der Frauen wird dafür nicht gebraucht. \quad f) Aussage 1 richtig: $10 + 45 = 55\,\%$, mehr als die Hälfte. Aussage 2 falsch: jeder 10. wäre $= 10\,\%$, im Diagramm sind es $30\,\%$; entscheidbar auch ohne die Anzahl der Befragten.}
